\documentclass[12pt]{article}

\usepackage[utf8]{inputenc}
\usepackage[T1]{fontenc}
\usepackage[brazil]{babel}
\usepackage{amsmath, amssymb}
\usepackage{tikz}
\usetikzlibrary{automata, positioning, arrows.meta}
\usepackage{listings}
\usepackage{xcolor}
\usepackage[a4paper, margin=2.5cm]{geometry}

\lstset{
  language=Python,
  basicstyle=\ttfamily\small,
  keywordstyle=\color{blue},
  commentstyle=\color{gray},
  numbers=left,
  frame=single,
  breaklines=true,
  inputencoding=utf8,
  extendedchars=true,
  literate={ã}{{\~a}}1 {á}{{\'a}}1 {é}{{\'e}}1 {ç}{{\c c}}1 {õ}{{\~o}}1 {ê}{{\^e}}1 {í}{{\'i}}1 {ó}{{\'o}}1
}

\title{Validador de URLs e Endereços IP: \\ Modelagem com Autômatos Finitos}
\author{Pedro \and Gus \and Guilherme \and Cauã}
\date{\today}

\begin{document}
\maketitle

\begin{abstract}
Este trabalho apresenta a modelagem de autômatos finitos determinísticos
para validação de URLs e endereços IP.
\end{abstract}

\section{Introdução}

\section{Validador de IP}
\subsection{Definição Formal}
% parte matemática:

\subsection{Diagrama de Transições}

\subsection{Implementação}

\section{Validador de URL}
\subsection{Definição Formal}
\subsection{Diagrama de Transições}
\subsection{Implementação}

\section{Conclusão}

\end{document}